%% ma: L11-B14-0019
%% lop: 11
%% chuong: 4
%% bai: 14
%% dang: 11.14.1
%% loai: TN4
%% muc: TH
%% dap_an: "B"
%% nguon: "tai-lieu-hoc-tap-mon-toan-lop-11-theo-chuong-trinh-giao-duc-pho-thong-2018.pdf (Lớp Toán Thầy Duy, Phạm Lê Duy, THPT Chu Văn An), Chương 4 Bài 5, Phần 1, Câu 14 (tr.529)"
%% kien_thuc: "Lớp 11 Bài 14 (phép chiếu song song)"
%% trang_thai: chinh-thuc
%% ghi_chu: "Đã duyệt 03/11/2026."
\begin{ex}
Cho hình hộp $ABCD.A^{\prime}B^{\prime}C^{\prime}D^{\prime}$. Hình biểu diễn của $\triangle A^{\prime}C^{\prime}D^{\prime}$ qua phép chiếu song song lên mặt phẳng $(ABCD)$ theo phương $A^{\prime}B$ là
\choice{Một tam giác $ABC$.}{\True Một tam giác $BCI$, với $I$ là điểm sao cho $C$ là trung điểm của $DI$.}{Một tam giác $ACD$.}{Một tam giác $DCM$, với $M$ là điểm sao cho $C$ là trung điểm của $BM$.}
\loigiai{\sloppy Theo phương $A^{\prime}B$: $A^{\prime}\mapsto B$. Ta có $D^{\prime}C\parallel A^{\prime}B$ ($A^{\prime}BCD^{\prime}$ là hình bình hành) nên $D^{\prime}\mapsto C$. Gọi $I$ là điểm đối xứng của $D$ qua $C$; khi đó $C^{\prime}I\parallel A^{\prime}B$ (vì $\overrightarrow{C^{\prime}I}=\overrightarrow{A^{\prime}B}$) nên $C^{\prime}\mapsto I$. Hình chiếu của $\triangle A^{\prime}C^{\prime}D^{\prime}$ là $\triangle BIC$ với $C$ là trung điểm $DI$.}
\end{ex}
